\chapter{Combinaciones lineales e independencia}
\label{cap:combinaciones-independencia}

Este capítulo estudia qué vectores pueden construirse a partir de otros y cuándo una representación lineal no contiene redundancias. Estos conceptos preparan la definición de base y dimensión \cite{lang1986,hoffmankunze1971}.

\section{Combinaciones lineales}
Suma finita de múltiplos escalares y lectura de expresiones lineales.

\section{Envolvente lineal y conjuntos generadores}
Subespacio generado, conjuntos generadores y descripción de span.

\section{Dependencia e independencia lineal}
Relaciones lineales y definición de independencia.

\section{Criterios de independencia}
Pruebas mediante ecuaciones homogéneas, pivotes y propiedades estructurales.

\section{Conjuntos generadores mínimos}
Redundancia, eliminación de vectores y caracterización de conjuntos mínimos.

\section{Relaciones lineales entre vectores}
Espacio de relaciones y lectura de dependencias entre generadores.
